\subsection{The Heap Subsystem and Object Allocation}

\subsubsection{The Managed Runtime Allocation Arena: PyMalloc and Private Heap Segmentation Subsystems}

When Python code creates an object, it does not normally ask the operating system directly for memory. The programmer writes:

\begin{verbatim}
x = [1, 2, 3]
\end{verbatim}

but CPython must allocate several runtime objects behind that line: at least a list object and references to the element objects. This allocation happens inside CPython's managed memory system, not through a direct user-visible call to \texttt{malloc()}.

The official CPython model describes Python memory management as involving a private heap containing Python objects and internal data structures. This private heap is managed by the Python memory manager, not by ordinary Python code. The user manipulates objects, but CPython manages the memory blocks behind those objects. :contentReference[oaicite:0]{index=0}

This distinction is essential:

\begin{quote}
Python objects live in a runtime-managed heap, not in programmer-controlled C-style storage.
\end{quote}

In C, the programmer may write:

\begin{verbatim}
int *p = malloc(10 * sizeof(int));
free(p);
\end{verbatim}

In Python, ordinary code does not allocate and release object memory this way. CPython decides where an object is placed, how much memory it needs, how the allocation is grouped, and when the memory can be reused.

\subsubsection*{Why CPython Uses Its Own Allocator}

CPython could theoretically ask the system allocator for every small object. But Python programs create many small objects: integers, tuples, strings, dictionaries, frames, iterators, temporary values, function objects, method objects, and countless short-lived internal structures.

Calling the platform allocator for every small object would be expensive. It would also produce fragmentation and poor locality. CPython therefore uses specialized allocation machinery for Python objects and internal buffers.

The Python memory manager has several layers. At the lowest level, a raw memory allocator obtains memory from the operating system. Above that, CPython has object-specific and size-specific allocators that manage smaller blocks more efficiently. The official documentation explicitly notes that different components handle sharing, segmentation, preallocation, caching, and object-specific allocation policies. :contentReference[oaicite:1]{index=1}

The important teaching idea is simple: CPython does not treat every object allocation as an isolated operating-system allocation. It reserves larger regions and subdivides them for Python's object workload.

\subsubsection*{Arenas, Pools, and Blocks}

The small-object allocator traditionally known as \texttt{pymalloc} organizes memory into layers. The exact constants may vary across versions and platforms, but the conceptual structure is stable enough for understanding:

\begin{verbatim}
arena
    large region obtained from the system allocator

pool
    subdivision inside an arena, usually dedicated to one size class

block
    small allocation unit returned for one object or internal structure
\end{verbatim}

A simplified diagram is:

\begin{verbatim}
CPython private heap
    arena
        pool for 32-byte blocks
            block
            block
            block

        pool for 64-byte blocks
            block
            block
            block

        pool for 128-byte blocks
            block
            block
            block
\end{verbatim}

The goal is to make small allocations fast. If CPython needs a small block, it can often take one from an existing pool instead of asking the operating system for new memory.

This is why Python object allocation is not the same as a direct call to C \texttt{malloc()} for every object. There may be a system allocation underneath at some point, but CPython often reuses already reserved memory managed inside its own allocator.

\subsubsection*{Private Heap Does Not Mean Separate Physical RAM}

The phrase ``private heap'' can mislead if read too literally. It does not mean Python owns a physically separate memory device. CPython is still a native process. Its memory ultimately comes from the operating system through normal virtual memory mechanisms.

The phrase means that CPython manages a region of process memory according to Python's object-allocation rules. The operating system sees process memory. CPython sees arenas, pools, blocks, object headers, type-specific layouts, free lists, and garbage-collection metadata.

So the layers are:

\begin{verbatim}
operating system virtual memory
    gives memory to CPython process

CPython memory manager
    organizes memory into allocator structures

Python object system
    places objects into allocated blocks

Python code
    manipulates references to those objects
\end{verbatim}

The programmer does not normally see these layers, but their effects appear in performance, memory usage, object lifetime, and extension-module rules.

\subsubsection*{Small Objects and Large Objects}

\texttt{pymalloc} is mainly useful for small objects and small internal allocations. Large allocations may be handled differently and may go closer to the platform allocator. This is practical: a tiny tuple, a small integer object, and a large byte buffer do not have the same allocation needs.

Small objects benefit from size-class pooling. Large objects are less suitable for tiny fixed-size pools and may need larger contiguous regions. CPython's memory manager therefore does not use one identical strategy for everything.

This matters for intuition. Python code creates objects uniformly at the source level, but CPython may allocate them through different paths internally. A small tuple, a dictionary table, a huge string, and a large list backing array are not all handled identically.

\subsubsection*{Object-Specific Allocation}

Above the general allocator, some object types have specialized allocation behavior. The official documentation notes that object-specific allocators operate on the same private heap and may use policies adapted to each object type. For example, integers, strings, tuples, and dictionaries have different storage requirements and speed-space tradeoffs. :contentReference[oaicite:2]{index=2}

This is why ``object allocation'' in Python is not one single action. Creating a dictionary is not the same internal operation as creating an integer. Creating a list involves a list object plus an internal array of object references. Creating a string involves character storage and string metadata. Creating a tuple involves fixed-size reference storage.

At the Python level, all of these are objects. At the CPython level, each type has its own layout and allocation needs.

\subsubsection*{Why Extension Modules Must Respect the Python Memory Manager}

This allocator structure is also why C extension modules must be careful. A native extension that creates Python objects must use the Python memory and object APIs, not randomly mix allocation and deallocation rules.

The official documentation warns that extension writers should not operate on Python objects with ordinary C library \texttt{malloc()}, \texttt{calloc()}, \texttt{realloc()}, and \texttt{free()} because mixing the C allocator and the Python memory manager can corrupt memory. :contentReference[oaicite:3]{index=3}

The reason is straightforward: the allocator that gives memory must be compatible with the allocator that releases it. If CPython allocates a Python object through its own allocator, freeing it through the wrong mechanism breaks the allocator's internal bookkeeping.

This is not a Python syntax issue. It is a concrete runtime ownership issue.

\subsubsection*{The Practical Mental Model}

The practical model is:

\begin{verbatim}
Python code asks for objects indirectly
    x = SomeObject()

CPython type machinery decides object layout
    how large?
    which type?
    which allocator path?

CPython memory manager provides memory
    often from private heap structures
    often through pymalloc for small blocks

object is initialized
    header
    type pointer
    value-specific fields
    possible GC metadata

Python name receives reference
    x -> object
\end{verbatim}

The object is not created in a vague abstract space. It is allocated inside the CPython process, usually through CPython's managed memory machinery. The name in Python receives a reference to that object. The object then remains alive as long as references or runtime structures keep it reachable.

This prepares the next allocation topic. Once the allocator can provide memory, CPython must actually instantiate objects: choose a type, allocate the right layout, initialize fields, connect references, and return a usable Python object to the executing frame.


\subsubsection{Explicit Dynamic Instantiation: Analyzing Instance Generation on the Heap}

Object creation in Python looks simple at the source level:

\begin{verbatim}
p = Point(2, 3)
\end{verbatim}

but this line hides several runtime steps. CPython must find the object bound to the name \texttt{Point}, determine that it is callable, execute the type's construction protocol, allocate memory for a new instance, initialize that instance, and finally bind the name \texttt{p} to the resulting object.

The important point is that Python objects are created dynamically. The program does not reserve a fixed compile-time storage slot for every future object. Instead, objects are produced while the program runs, placed in CPython-managed heap memory, and connected to the executing program through references.

A simplified model is:

\begin{verbatim}
source code
    p = Point(2, 3)

runtime behavior
    resolve name Point
    call class object Point
    allocate new instance object on heap
    initialize instance state
    bind name p to the new object
\end{verbatim}

This differs sharply from a C local variable such as:

\begin{verbatim}
struct Point p;
\end{verbatim}

In C, this may reserve storage directly in the current stack frame, depending on where the declaration appears. In Python, a statement such as \texttt{p = Point(2, 3)} does not create an object inside the variable \texttt{p}. It creates a heap object and binds the name \texttt{p} to it.

\subsubsection*{Calling a Class Means Running the Construction Protocol}

A class in Python is itself an object. When the programmer writes:

\begin{verbatim}
Point(2, 3)
\end{verbatim}

the class object \texttt{Point} is called. This call is not just a jump to one fixed constructor function. It follows Python's object-construction protocol.

At a simplified level, instance creation involves two conceptual phases:

\begin{itemize}
    \item allocation: create a new object of the requested type;
    \item initialization: set up the object's initial state.
\end{itemize}

In Python terms, these phases correspond roughly to:

\begin{verbatim}
__new__   -> creates and returns the new object
__init__  -> initializes the returned object
\end{verbatim}

For example:

\begin{verbatim}
class Point:
    def __init__(self, x, y):
        self.x = x
        self.y = y
\end{verbatim}

When the program executes:

\begin{verbatim}
p = Point(2, 3)
\end{verbatim}

CPython creates a new \texttt{Point} instance, then calls \texttt{\_\_init\_\_} with that instance as \texttt{self}. The assignments inside \texttt{\_\_init\_\_} attach state to the new object.

Conceptually:

\begin{verbatim}
new Point instance
    initially allocated as an object

__init__(self, 2, 3)
    self.x = 2
    self.y = 3

final object
    x -> int object 2
    y -> int object 3
\end{verbatim}

The object is created before \texttt{\_\_init\_\_} finishes. \texttt{\_\_init\_\_} does not return the object. It configures an object that has already been allocated.

\subsubsection*{The Instance Lives on the Managed Heap}

The new instance is placed in CPython-managed heap memory. The Python name receives a reference to it:

\begin{verbatim}
p  --->  Point instance object
            x -> int object 2
            y -> int object 3
\end{verbatim}

This is the central memory distinction:

\begin{quote}
The name is not the object. The name points to the object.
\end{quote}

If another name is assigned to the same object:

\begin{verbatim}
q = p
\end{verbatim}

then no new \texttt{Point} instance is created. The binding becomes:

\begin{verbatim}
p ----+
      |
      v
   Point instance object

q ----+
\end{verbatim}

Both names refer to the same heap object. If the object is mutable and one reference modifies it, the change is visible through the other reference.

For example:

\begin{verbatim}
q.x = 100
print(p.x)
\end{verbatim}

prints the modified value because \texttt{p} and \texttt{q} are not two independent objects. They are two references to one object.

\subsubsection*{Instance State Is Usually Stored in an Attribute Namespace}

Most ordinary Python objects store instance attributes in an instance namespace, commonly exposed as \texttt{\_\_dict\_\_}.

For example:

\begin{verbatim}
p = Point(2, 3)
print(p.__dict__)
\end{verbatim}

may show a mapping similar to:

\begin{verbatim}
{'x': 2, 'y': 3}
\end{verbatim}

Conceptually:

\begin{verbatim}
Point instance
    type -> Point
    instance namespace
        "x" -> int object 2
        "y" -> int object 3
\end{verbatim}

So the attributes \texttt{p.x} and \texttt{p.y} are not C struct fields in the simple fixed-layout sense. For ordinary dynamic Python objects, they are name-to-object bindings stored in an attribute namespace.

This is why Python objects can often receive new attributes after creation:

\begin{verbatim}
p.z = 9
\end{verbatim}

The instance namespace becomes:

\begin{verbatim}
"x" -> int object 2
"y" -> int object 3
"z" -> int object 9
\end{verbatim}

This flexibility has a cost. A dynamic attribute dictionary is heavier than a fixed C struct layout. Python gains dynamic object shape; CPython pays with extra memory and lookup machinery.

\subsubsection*{Slots Reduce Dynamic Attribute Storage}

Python also allows more fixed instance layouts through \texttt{\_\_slots\_\_}:

\begin{verbatim}
class Point:
    __slots__ = ("x", "y")

    def __init__(self, x, y):
        self.x = x
        self.y = y
\end{verbatim}

With slots, instances do not need the same ordinary per-object attribute dictionary for those attributes. CPython can store references in a more fixed internal layout.

The conceptual difference is:

\begin{verbatim}
ordinary instance
    attribute dictionary:
        "x" -> object
        "y" -> object

slotted instance
    fixed slots:
        slot x -> object
        slot y -> object
\end{verbatim}

This does not turn Python into C, but it shows that Python object layout can vary depending on class definition. The runtime still stores references to Python objects, but it may store them through a more compact structure.

\subsubsection*{Built-in Objects Use Type-Specific Layouts}

User-defined instances are only one case. Built-in objects use specialized internal layouts.

A list object is not a dictionary of attributes holding its elements. It is a list object with internal storage for references to element objects. A tuple has fixed reference storage. A dictionary has a hash-table structure. A string has a Unicode representation. An integer has a variable-sized numeric representation.

At source level:

\begin{verbatim}
items = [1, 2, 3]
\end{verbatim}

Conceptually:

\begin{verbatim}
items -> list object
            element slot 0 -> int object 1
            element slot 1 -> int object 2
            element slot 2 -> int object 3
\end{verbatim}

The list object itself lives on the managed heap. Its internal element storage contains references to other objects. The integers are also Python objects, though CPython may cache or reuse some common integer objects.

This means that container creation is object-graph creation. The runtime is not merely placing raw values in a flat memory region. It is building a network of objects and references.

\subsubsection*{Construction May Reuse Existing Objects}

Dynamic instantiation does not always mean a fresh object is created in the intuitive sense. CPython and Python types may reuse objects in some cases.

For example, small integers and some strings may be cached or interned. Immutable objects may be reused because their visible value cannot be changed. A constructor may also return an existing object from \texttt{\_\_new\_\_} instead of allocating a fresh one.

This means the programmer should not reason as follows:

\begin{quote}
Every expression that looks like construction necessarily creates a brand-new object.
\end{quote}

The safer model is:

\begin{quote}
An expression produces a reference to an object. Whether that object is newly allocated, cached, reused, or internally shared depends on the type and implementation.
\end{quote}

For mutable user-defined instances, a class call usually creates a fresh object. For immutable or optimized built-in objects, reuse is possible.

\subsubsection*{Instantiation Creates References, Not Deep Copies}

Object creation often connects existing objects into a new object. It does not automatically copy everything deeply.

Example:

\begin{verbatim}
data = [1, 2, 3]

class Box:
    def __init__(self, value):
        self.value = value

b = Box(data)
\end{verbatim}

The new \texttt{Box} instance does not contain an independent deep copy of \texttt{data}. It stores a reference to the same list object:

\begin{verbatim}
b -> Box instance
        value ----+
                  |
                  v
data ---------> list object [1, 2, 3]
\end{verbatim}

If the list is mutated:

\begin{verbatim}
data.append(4)
\end{verbatim}

then \texttt{b.value} observes the same changed list.

This is the same reference model as assignment, now appearing inside object construction. Initializing an instance attribute usually stores a reference to an object, not a deep structural clone.

\subsubsection*{The Practical Mental Model}

The practical model of instance generation is:

\begin{verbatim}
class call
    resolves class object

type machinery
    invokes allocation protocol

allocator
    reserves memory for a new object layout

object initialization
    fills object header and type-specific fields

__init__
    binds initial attribute references

assignment
    binds a Python name to the resulting object
\end{verbatim}

At the source level, this may be one line:

\begin{verbatim}
p = Point(2, 3)
\end{verbatim}

At the runtime level, it is a sequence of object operations:

\begin{verbatim}
resolve Point
call Point
allocate Point instance
initialize attributes
return object reference
bind p to that reference
\end{verbatim}

This explains why Python object creation is flexible but not free. Every object must have a type, a lifetime, reference relationships, and memory storage. Every dynamic attribute adds namespace structure. Every container connects references to other objects. CPython hides this machinery from the source code, but the machinery exists.

For an effective Python programmer, the essential conclusion is this:

\begin{quote}
Instantiation creates heap objects and reference relationships. Names and attributes do not contain objects directly; they point to objects managed by CPython.
\end{quote}


\subsubsection{Object Persistence Invariants: Reference Counting Foundations and Automated Slot Reclamation Realities}

Once a Python object has been allocated, CPython must decide how long that object remains alive. This is the problem of object persistence. An object should remain valid while the program can still reach it, and it should become reclaimable when the program can no longer use it. CPython solves most of this problem through reference counting, supplemented by cyclic garbage collection.

The central invariant is simple:

\begin{quote}
An object remains alive while there are active strong references to it.
\end{quote}

A reference may be held by a name, a container, an attribute, a frame, a closure cell, an exception object, an internal runtime structure, or a temporary value on the bytecode evaluation stack. The object does not care whether the reference came from a variable, a list, a dictionary, or a function call. From the memory-management point of view, what matters is that something still points to it.

\subsubsection*{References Keep Objects Alive}

Consider:

\begin{verbatim}
a = []
b = a
\end{verbatim}

Only one list object has been created. Both names refer to it:

\begin{verbatim}
a ----+
      |
      v
   list object
      ^
      |
b ----+
\end{verbatim}

The object persists because references to it exist. If one name is rebound:

\begin{verbatim}
a = None
\end{verbatim}

the list object is not automatically destroyed, because \texttt{b} still refers to it:

\begin{verbatim}
a -> None object

b -> list object
\end{verbatim}

Only when the last strong reference disappears can the object become reclaimable.

\subsubsection*{Reference Count Updates}

At the CPython level, an object carries a reference count. When a new strong reference is created, the count is increased. When a strong reference is released, the count is decreased.

Conceptually:

\begin{verbatim}
a = object()
    reference count becomes 1

b = a
    reference count becomes 2

a = None
    reference count becomes 1

b = None
    reference count becomes 0
    object can be deallocated
\end{verbatim}

This is not exact visible Python code behavior in every optimized internal case, but it is the correct conceptual model for CPython object lifetime.

Reference count changes happen constantly. Assignment, function calls, returns, container insertion, container removal, attribute assignment, frame creation, frame destruction, and exception handling can all create or release references.

For example:

\begin{verbatim}
items = []
items.append(object())
\end{verbatim}

The newly created object remains alive because the list now holds a reference to it. If the list is later cleared:

\begin{verbatim}
items.clear()
\end{verbatim}

the list releases its reference to the object. If no other reference exists, the object can be destroyed.

\subsubsection*{Slot Reclamation Means Releasing Stored References}

A container does not merely hold abstract values. It holds references to objects. When a container slot is overwritten or removed, CPython must release the old reference stored in that slot.

Example:

\begin{verbatim}
items = [object()]
items[0] = "new"
\end{verbatim}

Before replacement:

\begin{verbatim}
items -> list object
            slot 0 -> old object
\end{verbatim}

After replacement:

\begin{verbatim}
items -> list object
            slot 0 -> string object "new"
\end{verbatim}

The list slot no longer refers to the old object. CPython must decrease the old object's reference count. If that slot was the last strong reference, the old object can be deallocated.

This is what automated slot reclamation means in practice. When a name, attribute, list position, dictionary entry, tuple during destruction, frame local slot, or internal temporary slot stops holding a reference, CPython releases that reference. The memory discipline is attached to reference ownership, not to source-level variable declarations.

\subsubsection*{Name Rebinding Also Releases References}

The same idea applies to ordinary names:

\begin{verbatim}
x = []
x = 10
\end{verbatim}

The second assignment does not mutate the list into an integer. It rebinds the name \texttt{x}. The old reference from \texttt{x} to the list is released. The new reference from \texttt{x} points to the integer object.

Conceptually:

\begin{verbatim}
before:
    x -> list object

after:
    x -> int object 10
